
Foundation model pretraining typically uses static domain mixing ratios—fixed proportions maintained throughout training. This paper reports a proof-of-concept implementation of diversity-ranked domain scheduling, which orders training domains from high to low diversity using smooth Gaussian-weighted transitions. The implementation uses corpus-level statistics (vocabulary entropy, syntactic complexity, semantic spread) to rank six domains from The Pile dataset, yielding scores from 0.92 (Pile-CC) to 0.35 (PubMed Central). All core components execute correctly: 22 unit tests pass, the curriculum scheduler produces valid probability distributions with proper normalization constraints, and a smoke test (10 training steps, GPT-2 style model with 760M parameters) completes without errors. The smoke test composite benchmark score of 0.2558 reflects an untrained model and provides no evidence regarding performance improvements. The proposed hypothesis—that diversity-ranked scheduling improves model performance through path-dependent gradient covariance geometry—remains untested. Full-scale experiments (100,000 steps at 1B scale, 150,000 steps at 7B scale, n=5 seeds per condition) are required to test whether diversity-ranked scheduling outperforms static mixture baselines by ≥2.0\% at 1B scale and ≥0.5\% at 7B scale. The mechanistic explanation requires gradient covariance measurements (Participation Ratio, CKA similarity, Fisher overlap) not performed in this work. This proof-of-concept establishes implementation feasibility prior to committing computational resources (estimated 45,000 GPU-hours) to full validation.
